
\documentclass[11pt]{article}

\usepackage[margin=1.08in]{geometry}
\usepackage{amsmath,amssymb,amsthm,mathtools}
\usepackage{enumitem}
\usepackage{microtype}
\usepackage{hyperref}
\usepackage{booktabs}
\usepackage{xcolor}

\hypersetup{
  colorlinks=true,
  linkcolor=blue,
  citecolor=blue,
  urlcolor=blue
}

\newtheorem{theorem}{Theorem}[section]
\newtheorem{proposition}[theorem]{Proposition}
\newtheorem{lemma}[theorem]{Lemma}
\newtheorem{corollary}[theorem]{Corollary}
\newtheorem{definition}[theorem]{Definition}
\newtheorem{remark}[theorem]{Remark}
\newtheorem{assumption}[theorem]{Assumption}

\newcommand{\R}{\mathbb R}
\newcommand{\T}{\mathbb T}
\newcommand{\cM}{\mathcal M}
\newcommand{\cR}{\mathcal R}
\newcommand{\cN}{\mathcal N}
\newcommand{\cO}{\mathcal O}
\newcommand{\Wu}{W^{u}}
\newcommand{\Ws}{W^{s}}
\newcommand{\Fix}{\operatorname{Fix}}
\newcommand{\htop}{h_{\mathrm{top}}}
\newcommand{\Span}{\operatorname{span}}
\newcommand{\pr}{p_r}
\newcommand{\pz}{p_z}
\newcommand{\eps}{\varepsilon}

\title{\bf A rigorous computer-assisted scheme for positive topological entropy\\
in the spatial isosceles three-body problem}

\author{
Lei Liu\thanks{School of Mathematics, Shandong University, Jinan, China.
E-mail: \texttt{liulei30@email.sdu.edu.cn}}
\and
Pedro A. S. Salom\~ao\thanks{Shenzhen International Center for Mathematics,
Southern University of Science and Technology (SUSTech), Shenzhen, China.
E-mail: \texttt{psalomao@sustech.edu.cn}}
}

\date{}

\begin{document}
\maketitle

\begin{abstract}
A problem of Jean-Pierre Marco, recorded as Problem 6 in the list of
Albouy--Cabral--Santos, asks whether the collisionless isosceles three-body
problem has positive topological entropy whenever the angular momentum and
the energy determine a smooth compact integral manifold.
We develop a rigorous computer-assisted scheme for a concrete compact energy
surface of the spatial isosceles three-body problem.  After fixing the energy
to $-1$, we consider
\[
 \beta=\frac{3}{100},\qquad e=\frac{987}{1000},
\]
for which $\beta^2+e^2<1$.
The proof is reduced to a finite list of interval inequalities for a
Poincar\'e map.  Once validated, these inequalities imply the existence of a
symmetric hyperbolic period-four orbit, a local unstable graph, and a
symmetric homoclinic topological crossing.  The geometric part of the argument is independent of
the numerical implementation: reversibility turns the zero of a scalar
apoapsis coordinate into a homoclinic point, while a one-sided derivative
bound forces a topological crossing.  The theorem of Burns and Weiss then
gives positive topological entropy.  A stronger derivative certificate gives
a transverse homoclinic orbit and hence a Smale horseshoe.
We explain in detail how interval arithmetic enters the argument and isolate
the exact finite validation statement needed for a complete
computer-assisted proof.
\end{abstract}

\medskip
\noindent
\textbf{2020 Mathematics Subject Classification.}
Primary 37J46, 37J55; Secondary 70F07, 37B40, 65G20.

\medskip
\noindent
\textbf{Keywords.}
spatial isosceles three-body problem, topological entropy, homoclinic orbit,
reversible dynamics, interval arithmetic, CAPD, Poincar\'e map.

\section{Introduction}

The isosceles three-body problem is one of the classical invariant
subsystems of the Newtonian three-body problem.  Two bodies have equal mass
and remain symmetric with respect to a fixed axis, while the third body moves
on that axis.  After reduction by the rotational symmetry, the system has two
degrees of freedom.  It is simple enough to admit detailed geometric and
analytic study, but rich enough to display many of the mechanisms that occur
in general celestial mechanics.

A basic question is whether the dynamics on compact regular energy surfaces
is genuinely chaotic.  In the list of open problems compiled by
Albouy--Cabral--Santos \cite{AlbouyCabralSantos}, Problem 6 is attributed to
Jean-Pierre Marco and asks whether the topological entropy is positive in
the collisionless isosceles three-body problem for values of the angular
momentum and energy for which the integral manifold is a smooth compact
manifold.  We refer to this question as the \emph{Marco entropy problem}.
The Hamiltonian used below is obtained after fixing the angular momentum and
reducing the cyclic rotation angle.  If one formulates Marco's problem on the
unreduced integral manifold, the quotient map onto the reduced energy surface
semiconjugates the two flows.  Since topological entropy cannot increase under
a factor map, positive entropy for the reduced flow implies positive entropy
for the corresponding unreduced flow.  Thus a positive result below gives a
partial affirmative answer to Marco's problem in either formulation.

There are classical constructions of symbolic dynamics in the isosceles
problem.  In particular, Moeckel \cite{Moeckel1984} constructed invariant
sets with symbolic dynamics and heteroclinic phenomena in regimes related to
close approaches to triple collision.  These results are fundamental, but
they do not settle the general compact problem formulated by Marco.

More recently, symplectic and contact methods have provided a rather precise
description of the compact spatial isosceles problem.  For suitable
parameters the compact energy surface is diffeomorphic to $S^3$, and the
Euler orbit binds disk-like global surfaces of section; see
\cite{HuLiuOuSalomaoYu}.  Quantitative constraints coming from embedded
contact homology imply, in particular, the existence of infinitely many
periodic orbits on every compact energy surface in the regime considered in
\cite{ECHconstraints}.  Nevertheless, infinitely many periodic orbits do not
by themselves force positive topological entropy.  Even an integrable twist
map may possess infinitely many periodic orbits and have zero entropy.

The aim of the present paper is different.  We look for a concrete
hyperbolic mechanism: a homoclinic crossing associated with a symmetric
hyperbolic periodic orbit.  The existence of such a crossing can be reduced
to a finite number of inequalities for a two-dimensional Poincar\'e map.
This makes the problem particularly well suited to interval arithmetic.

The parameter values used throughout are
\begin{equation}\label{eq:parameters}
 \beta=\frac{3}{100},\qquad
 e=\frac{987}{1000},\qquad
 H=-1.
\end{equation}
They correspond to
\begin{equation}\label{eq:alphaomega}
 \alpha=\frac{12}{97},\qquad
 \varpi^2=\frac{25831}{1800}.
\end{equation}
Moreover,
\[
 \beta^2+e^2=0.975069<1,
\]
so the energy surface lies strictly in the compact regime.

The computer-assisted part of the argument is deliberately kept small.  It
has four roles:
\begin{enumerate}[label=(\roman*)]
\item validate a symmetric period-four point of a periapsis Poincar\'e map;
\item prove that this point is hyperbolic;
\item prove cone, trapping and exit inequalities giving the true local
unstable manifold as a graph;
\item propagate a small part of this graph to an apoapsis section and prove a
sign change together with a derivative inequality.
\end{enumerate}
Everything after these inequalities is an exact geometric argument.

The main mathematical reduction is the following.

\begin{theorem}[Finite entropy certificate]\label{thm:certificate}
Consider the spatial isosceles three-body problem with the parameters
\eqref{eq:parameters}.  Assume that the interval inequalities
\textnormal{(C1)--(C4)} of Section~\ref{sec:certificate} are validated.
Then the Hamiltonian flow on $\cM=H^{-1}(-1)$ has positive topological
entropy:
\[
 \htop(\phi_H|_{\cM})>0.
\]
More precisely, the flow has a hyperbolic periodic orbit with a homoclinic
topological crossing.  If the stronger transverse version
\textnormal{(C4$^+$)} is validated, then the stable and unstable manifolds
have a transverse homoclinic intersection and the flow contains a suspended
Smale horseshoe.
\end{theorem}

The significance of Theorem~\ref{thm:certificate} is that it reduces the
remaining computer-assisted step to finite interval statements involving only explicit formulas for the vector
field, the Poincar\'e maps and their first derivatives.  No numerical
unstable-manifold polynomial is assumed.  The local unstable manifold is
obtained from a rigorous cone argument.

\begin{remark}[Status of the computer-assisted step]\label{rem:status}
The theorem above separates the mathematical proof from the validation run.
A final computer-assisted proof of positive entropy at
\eqref{eq:parameters} consists of Theorem~\ref{thm:certificate} together
with a machine-verifiable output certifying \textnormal{(C1)--(C4)}.
Exploratory floating-point computations locate the relevant periodic point
and homoclinic crossing with large margins; see
Section~\ref{sec:exploratory}.  These floating-point values are used only to
center the interval calculation and are not part of the proof.
\end{remark}

The paper is organized as follows.  Section~\ref{sec:model} recalls the
Hamiltonian and the compact parameter regime.  Section~\ref{sec:symmetry}
describes the reversible Poincar\'e map and the symmetric period-four orbit.
Section~\ref{sec:graph} gives a complete proof of the local unstable graph
lemma.  Section~\ref{sec:homoclinic} turns an apoapsis sign change into a
homoclinic point.  Section~\ref{sec:crossing} proves that one derivative
bound is enough to obtain a topological crossing and positive entropy.
Section~\ref{sec:interval} explains the use of interval arithmetic in the
problem.  Section~\ref{sec:certificate} gives the precise finite validation
statement.

\section{The spatial isosceles three-body problem}\label{sec:model}

We use coordinates
\[
 (r,z,\pr,\pz)\in \R_+\times\R\times\R^2.
\]
The reduced Hamiltonian is
\begin{equation}\label{eq:H}
 H(r,z,\pr,\pz)
 =
 \frac{\pr^2+\pz^2}{2}
 +\frac{\varpi^2}{2r^2}
 -\frac1r
 -\frac{4}{\alpha\sqrt{r^2+(1+2\alpha)z^2}}.
\end{equation}
Hamilton's equations are
\begin{equation}\label{eq:HamEq}
\begin{aligned}
 \dot r&=\pr,\\
 \dot z&=\pz,\\
 \dot{\pr}
 &=
 \frac{\varpi^2}{r^3}
 -\frac1{r^2}
 -\frac{4r}
 {\alpha\bigl(r^2+(1+2\alpha)z^2\bigr)^{3/2}},\\
 \dot{\pz}
 &=
 -\frac{4(1+2\alpha)z}
 {\alpha\bigl(r^2+(1+2\alpha)z^2\bigr)^{3/2}}.
\end{aligned}
\end{equation}

Following the standard parametrization, set
\begin{equation}\label{eq:betae}
 \beta=\frac{1}{1+4/\alpha},
 \qquad
 e=\sqrt{1+2h\varpi^2\beta^2}.
\end{equation}
For negative energy we may use the standard scaling to fix $h=-1$.
The scaling conjugates the phase portraits and changes time only by a
positive constant factor, so it does not affect whether the topological
entropy is positive.  Hence
\[
 \varpi^2=\frac{1-e^2}{2\beta^2}.
\]
For \eqref{eq:parameters} this gives exactly \eqref{eq:alphaomega}.

For
\[
 0<\beta^2+e^2<1,
\]
the regular energy hypersurface $\cM=H^{-1}(-1)$ is compact and
diffeomorphic to $S^3$; see \cite{Moeckel1984,HuLiuOuSalomaoYu}.
The centrifugal term prevents $r$ from approaching zero, so the flow on
$\cM$ is collisionless and complete.

\subsection{The energy graph on the periapsis section}

The local Poincar\'e sections used below lie in $\{\pr=0\}$.  Near the
periodic orbit of interest, the energy equation can be solved uniquely for
$r$ as a function of $(\pz,z)$.

Set
\[
 \Phi(r;p,z):=H(r,z,0,p)+1.
\]
Then
\begin{equation}\label{eq:Phi_r}
 \partial_r\Phi
 =
 -\frac{\varpi^2}{r^3}
 +\frac{1}{r^2}
 +\frac{4r}
 {\alpha\bigl(r^2+(1+2\alpha)z^2\bigr)^{3/2}}.
\end{equation}

\begin{lemma}[Validated energy graph]\label{lem:energygraph}
For the parameters \eqref{eq:parameters}, if
\[
 8.5767\le p\le8.5770,\qquad |z|\le10^{-4},
\]
then the equation $\Phi(r;p,z)=0$ has a unique solution
\[
 r=r(p,z)\in(0.36,0.39).
\]
Moreover $r(p,z)$ is real analytic.
\end{lemma}

\begin{proof}
Since
\[
 \bigl(r^2+(1+2\alpha)z^2\bigr)^{3/2}\ge r^3,
\]
we have
\[
 \partial_r\Phi
 \le
 -\frac{\varpi^2}{r^3}
 +\frac{1+4/\alpha}{r^2}
 =
 \frac{r/\beta-\varpi^2}{r^3}.
\]
A direct interval evaluation on $0.36\le r\le0.39$ gives
\begin{equation}\label{eq:Hrbound}
 \partial_r\Phi<-22.7.
\end{equation}
At the two vertical sides one obtains the uniform bounds
\begin{equation}\label{eq:PhiSides}
 \Phi(0.36;p,z)>0.552,
 \qquad
 \Phi(0.39;p,z)<-0.512.
\end{equation}
Thus the intermediate value theorem gives existence and
\eqref{eq:Hrbound} gives uniqueness.  Analyticity follows from the analytic
implicit function theorem.
\end{proof}

\begin{remark}\label{rem:newton}
Lemma~\ref{lem:energygraph} is important for the computer-assisted proof.
An interval Newton iteration may be used to \emph{narrow} the enclosure of
$r(p,z)$, but existence must first be certified.  Monotonicity together with
the endpoint signs \eqref{eq:PhiSides} provides this certification.  Merely
checking that $\partial_r\Phi$ does not vanish would prove uniqueness but not
existence for every parameter pair in a box.
\end{remark}

Differentiating the identity $\Phi(r(p,z);p,z)=0$ gives
\begin{equation}\label{eq:liftderivatives}
 r_p=-\frac{p}{H_r},
 \qquad
 r_z=-\frac{H_z}{H_r}.
\end{equation}
These formulas are used to convert the derivative of the four-dimensional
flow to the derivative of the two-dimensional Poincar\'e map in coordinates
$(\pz,z)$.

\section{Reversibility and a symmetric period-four orbit}\label{sec:symmetry}

Let $\Sigma_{\rm per}$ be the periapsis section
\[
 \Sigma_{\rm per}
 =
 \{\pr=0,\ \dot{\pr}>0\}\cap\cM,
\]
and let $P$ denote the first return map from one periapsis crossing to the
next.  We use coordinates
\[
 (p,z)=(\pz,z)
\]
on $\Sigma_{\rm per}$.

The Hamiltonian is invariant under the symplectic involution
\[
 \mathcal C(r,z,\pr,\pz)
 =
 (r,-z,\pr,-\pz)
\]
and under the standard time reversal
\[
 \mathcal T(r,z,\pr,\pz)
 =
 (r,z,-\pr,-\pz).
\]
The composition
\begin{equation}\label{eq:fullreverser}
 \cR:=\mathcal C\mathcal T,\qquad
 \cR(r,z,\pr,\pz)=(r,-z,-\pr,\pz),
\end{equation}
is an anti-symplectic reverser.

On the periapsis section these maps induce
\[
 C(p,z)=(-p,-z),\qquad
 R(p,z)=(p,-z),\qquad
 S=CR=(-p,z).
\]
They satisfy
\begin{equation}\label{eq:reversibility}
 CP=PC,\qquad
 RPR=P^{-1},\qquad
 SPS=P^{-1}.
\end{equation}

For $q=(p,0)\in\Fix R$ define the scalar function
\begin{equation}\label{eq:g}
 g(p):=\pi_p P(p,0).
\end{equation}

\begin{lemma}[Symmetric period-four criterion]\label{lem:period4}
If $p_*\ne0$ and $g(p_*)=0$, then
\[
 q=(p_*,0)
\]
satisfies
\[
 P^2q=Cq,\qquad P^4q=q.
\]
In particular $q$ has prime period four for $P$.
\end{lemma}

\begin{proof}
Since $q\in\Fix R$ and $g(p_*)=0$, we have
$Pq\in\Fix S$.  Therefore
\[
 Pq=SPq=CRPq.
\]
Using $RP=P^{-1}R$, $Rq=q$ and $CP=PC$, we obtain
\[
 Pq=CP^{-1}q=P^{-1}Cq.
\]
Hence $P^2q=Cq$.  Applying $P^2$ once more and using commutation with $C$
gives
\[
 P^4q=P^2Cq=CP^2q=C^2q=q.
\]
Since $p_*\ne0$, $Cq\ne q$, hence $P^2q\ne q$; also $Pq$ belongs to
$\Fix S$ while $q\notin\Fix S$.  Thus the prime period is four.
\end{proof}

We set
\begin{equation}\label{eq:F}
 F:=P^4.
\end{equation}
The interval computation searches for $p_*$ in
\begin{equation}\label{eq:pbracket}
 I_*=[8.5768,8.5769].
\end{equation}
The validation criterion will prove a strict sign change of $g$ at the
endpoints and $g'>0$ on $I_*$.  It will also prove that the two multipliers
of $DF(q)$ satisfy
\begin{equation}\label{eq:multipliertarget}
 5.5<\lambda_u<5.8,
 \qquad
 0.17<\lambda_s<0.19.
\end{equation}
Consequently $q$ is a hyperbolic fixed point of $F$.

\section{A rigorous local unstable graph}\label{sec:graph}

We now prove the only local invariant-manifold statement needed later.
The point is to avoid treating a numerically computed Taylor polynomial as
the actual unstable manifold.

Let $B$ be the linear coordinate matrix
\begin{equation}\label{eq:B}
 B=
 \begin{pmatrix}
 0.68026252175126 & 0.68026252100350\\
 -0.73296855423723& 0.73296855493122
 \end{pmatrix},
\end{equation}
whose columns approximate unstable and stable eigenvectors.
Write
\[
 (p,z)=q+B(u,v).
\]
For $U>0$ and $a>0$ define
\begin{equation}\label{eq:N}
 N=[-U,U]\times[-aU,aU].
\end{equation}
Let
\[
 \widehat F(u,v)
 =
 B^{-1}\bigl(F(q+B(u,v))-q\bigr)
\]
denote $F$ in these coordinates and write
\[
 D\widehat F=
 \begin{pmatrix}
 J_{11}&J_{12}\\
 J_{21}&J_{22}
 \end{pmatrix}.
\]

\begin{proposition}[Cone--graph criterion]\label{prop:graph}
Assume that $\widehat F$ is defined on $N$, fixes the origin, and that the
following estimates hold on $N$:
\begin{align}
 \mu
 &:=
 \inf_N |J_{11}|
 -a\sup_N|J_{12}|>1,
 \label{eq:cone1}\\
 \sup_N|J_{21}|+a\sup_N|J_{22}|
 &<
 a\mu,
 \label{eq:cone2}\\
 |\pi_v\widehat F(u,v)|
 &<aU
 \qquad\text{for every }(u,v)\in N,
 \label{eq:trap}\\
 \pi_u\widehat F(U,v)&>U,
 \qquad
 \pi_u\widehat F(-U,v)<-U
 \quad (|v|\le aU).
 \label{eq:exit}
\end{align}
Then the local unstable manifold contains a graph
\begin{equation}\label{eq:unstablegraph}
 \Gamma^u
 =
 \{(u,h(u)):-U\le u\le U\}
 \subset \Wu_{\rm loc}(0;\widehat F)
\end{equation}
with
\begin{equation}\label{eq:hslope}
 h(0)=0,\qquad
 |h'(u)|\le a,\qquad
 |h(u)|\le a|u|.
\end{equation}
In particular, for every $u\in[-U,U]$ the point $(u,h(u))$ is a
genuine point of the unstable manifold.
\end{proposition}

\begin{proof}
Let $\mathcal G$ denote the set of Lipschitz functions
\[
 h:[-U,U]\longrightarrow[-aU,aU]
\]
such that $h(0)=0$ and $\operatorname{Lip}(h)\le a$.
This is a compact convex subset of $C^0([-U,U])$.

For $h\in\mathcal G$, consider
\[
 \Gamma_h=\{(u,h(u)):-U\le u\le U\}.
\]
The function
\[
 U_h(u):=\pi_u\widehat F(u,h(u))
\]
is absolutely continuous.  By \eqref{eq:cone1}, $J_{11}$ never vanishes on
the connected set $N$, hence it has constant sign there.  At almost every
point,
\[
 |U_h'(u)|
 =
 |J_{11}+J_{12}h'(u)|
 \ge
 \inf_N|J_{11}|-a\sup_N|J_{12}|
 =
 \mu>1.
\]
The sign of $U_h'$ is therefore constant.  The exit inequalities
\eqref{eq:exit} imply
\[
 U_h(-U)<-U<U<U_h(U),
\]
so $U_h$ cannot be decreasing.  Consequently
\[
 U_h'(u)\ge\mu>1
\]
almost everywhere.  Thus $U_h$ is strictly increasing, and the portion of
$\widehat F(\Gamma_h)$ lying over $[-U,U]$ is defined on the whole
horizontal interval.

At almost every point of $\Gamma_h$, the slope of the image is
\[
 \frac{J_{21}+J_{22}h'}
 {J_{11}+J_{12}h'}.
\]
By \eqref{eq:cone1}--\eqref{eq:cone2}, its absolute value is strictly less
than $a$.  The trapping estimate \eqref{eq:trap} shows that the image stays
strictly between the horizontal faces of $N$.  Therefore
\[
 \widehat F(\Gamma_h)\cap N
\]
is again a graph $\Gamma_{Th}$ for some $Th\in\mathcal G$.
Since $\widehat F(0)=0$, we have $(Th)(0)=0$.

The expansion estimate $\mu>1$ gives uniform control of the inverse of
$U_h$, and therefore the graph transform
\[
 T:\mathcal G\longrightarrow\mathcal G
\]
is continuous in the $C^0$ topology.  By Schauder's fixed point theorem
there exists $h\in\mathcal G$ with $Th=h$.

Let $x=(u,h(u))$ belong to the fixed graph.  There is a unique inverse
iterate $x_{-1}=(u_{-1},h(u_{-1}))$ in the same graph.  Since
$U_h(0)=0$ and $U_h'\ge\mu$ almost everywhere,
\[
 |u|\ge\mu|u_{-1}|.
\]
Iterating,
\[
 |u_{-n}|\le\mu^{-n}|u|.
\]
Because $|h(u_{-n})|\le a|u_{-n}|$, we obtain
\[
 \widehat F^{-n}(x)\longrightarrow0.
\]
Hence the invariant graph is contained in $\Wu(0;\widehat F)$.

Since $h(0)=0$ and $\operatorname{Lip}(h)\le a$,
\[
 |h(u)|\le a|u|.
\]
The system and $F$ are real analytic in the collisionless region.
The graph is a connected interval through the hyperbolic fixed point and is
contained in its one-dimensional analytic unstable manifold.  Since the
$u$-projection is one-to-one on the graph, the analytic unstable-manifold
theorem implies that $h$ is real analytic in the interior of its interval.
\end{proof}

\begin{remark}
The proof uses only interval bounds on $D\widehat F$ and on the images of
the three relevant sets: the whole block and the two vertical faces.
No high-order parametrization of $\Wu(q)$ is required.
\end{remark}

\section{From a sign change to a symmetric homoclinic orbit}
\label{sec:homoclinic}

Let
\[
 \Sigma_{\rm apo}
 =
 \{\pr=0,\ \dot{\pr}<0\}\cap\cM
\]
be the apoapsis section and let
\[
 A:\Sigma_{\rm per}\longrightarrow\Sigma_{\rm apo}
\]
denote the first hitting map.  Define
\begin{equation}\label{eq:G}
 G:=A\circ F^6=A\circ P^{24}.
\end{equation}

Let the unstable graph of Proposition~\ref{prop:graph} be written in
physical section coordinates as
\[
 x(u)=q+B(u,h(u)).
\]
We use
\begin{equation}\label{eq:uLR}
 u_L=2.68\times10^{-5},
 \qquad
 u_R=2.70\times10^{-5}.
\end{equation}
The bound \eqref{eq:hslope} implies
\[
 h(u_j)\in[-a|u_j|,a|u_j|],
 \qquad j=L,R.
\]
Hence vertical interval tubes at $u_L,u_R$ rigorously contain the true
unstable-manifold points.

Define the apoapsis arc
\begin{equation}\label{eq:gamma}
 \gamma(u):=G(x(u)),
 \qquad u\in[u_L,u_R].
\end{equation}

\begin{proposition}[Symmetric homoclinic criterion]\label{prop:homoclinic}
Assume that $G$ is well defined on the relevant unstable tubes, that the
apoapsis crossings are regular, and that
\begin{equation}\label{eq:zsign}
 z(\gamma(u_L))<0<z(\gamma(u_R)).
\end{equation}
Then the flow has a homoclinic orbit to the periodic orbit $\cO$ generated by
$q$.
\end{proposition}

\begin{proof}
By continuity there exists $u_*\in(u_L,u_R)$ such that
\[
 z(\gamma(u_*))=0.
\]
Since $\gamma(u_*)\in\Sigma_{\rm apo}$, also $\pr=0$ there.
Thus, by \eqref{eq:fullreverser},
\[
 \cR(\gamma(u_*))=\gamma(u_*).
\]

The periodic orbit $\cO$ is invariant under $\cR$, and $\cR$ reverses time.
Therefore
\[
 \cR(\Wu(\cO))=\Ws(\cO).
\]
On the other hand $\gamma(u_*)$ lies on $\Wu(\cO)$ because it is obtained by
forward propagation of a point of the unstable graph.  Since it is fixed by
$\cR$, it also belongs to $\Ws(\cO)$.  Hence
\[
 \gamma(u_*)\in\Wu(\cO)\cap\Ws(\cO).
\]
This intersection is nontrivial.  Indeed $u_*>u_L>0$, so
$x(u_*)\ne q$.  If $\gamma(u_*)$ belonged to the periodic orbit $\cO$, then
invariance of $\cO$ under the flow would imply $x(u_*)\in\cO$.  On the
periapsis section, $x(u_*)$ would then be a point fixed by $F=P^4$.  But
$x(u_*)\in\Wu(q;F)$, hence $F^{-n}(x(u_*))\to q$; a fixed point of $F$ with
this property must be $q$, a contradiction.  Thus the trajectory through
$\gamma(u_*)$ is a nontrivial homoclinic orbit to $\cO$.
\end{proof}

\section{Topological crossing and positive entropy}\label{sec:crossing}

A homoclinic orbit alone does not imply positive entropy: stable and unstable
manifolds could coincide along a separatrix connection.  We therefore need a
crossing condition.

The apoapsis section is invariant under the restriction of the reverser
\eqref{eq:fullreverser}; in coordinates $(p,z)=(\pz,z)$ this restriction is
simply
\begin{equation}\label{eq:apoReflection}
 (p,z)\longmapsto(p,-z).
\end{equation}

The following criterion is tailored to interval validation.

\begin{proposition}[One derivative is enough]\label{prop:crossing}
In addition to the assumptions of Proposition~\ref{prop:homoclinic}, assume
that
\begin{equation}\label{eq:pmonotone}
 \frac{d}{du}\,\pi_p\gamma(u)<-\eta<0
 \qquad\text{for every }u\in[u_L,u_R].
\end{equation}
Then at least one point of the homoclinic intersection obtained from
\eqref{eq:zsign} is a topological crossing between
$\Wu(\cO)$ and $\Ws(\cO)$.
\end{proposition}

\begin{proof}
Write
\[
 p(u)=\pi_p\gamma(u),\qquad
 z(u)=\pi_z\gamma(u).
\]
By \eqref{eq:pmonotone}, $p$ is a real-analytic local coordinate along the
unstable arc.  Thus the unstable manifold can be written near the relevant
part of the apoapsis section as
\[
 z=\varphi(p)
\]
for a real-analytic function $\varphi$.

The endpoint condition \eqref{eq:zsign} implies that $\varphi$ assumes both
signs.  Since $\varphi$ is analytic and is not identically zero, one may
choose a zero $p_0$ across which its sign changes.  This zero has odd finite
multiplicity:
\[
 \varphi(p)
 =
 c(p-p_0)^{2k+1}
 +O((p-p_0)^{2k+2}),
 \qquad c\ne0.
\]

By the reflection symmetry \eqref{eq:apoReflection}, the corresponding
stable arc is
\[
 z=-\varphi(p).
\]
Consequently the two arcs exchange their vertical order when $p$ passes
through $p_0$.  This is a topological crossing in the sense of
Burns--Weiss \cite{BurnsWeiss}.
\end{proof}

\begin{theorem}\label{thm:entropyfromcrossing}
Under the assumptions of Proposition~\ref{prop:crossing},
\[
 \htop(\phi_H|_{\cM})>0.
\]
\end{theorem}

\begin{proof}
Choose a disk-like global surface of section for the compact energy
surface; such sections are available in the compact SI3BP by
\cite{HuLiuOuSalomaoYu}.  In the present setting the first return map and the
return time admit smooth extensions to the closed disk; see
\cite[Propositions~4.1--4.2]{ECHconstraints}.  A local holonomy map obtained
by flowing from the apoapsis section to this disk is a diffeomorphism near
the homoclinic point and therefore preserves the topological crossing.

Let $\psi$ be the first return map of the disk.  Choose one intersection
$x_0\in\cO$ with the disk and let $k$ be its period under $\psi$.  After
choosing the appropriate intersection of the homoclinic orbit with the same
disk, we obtain a homoclinic topological crossing for the hyperbolic fixed
point $x_0$ of $\psi^k$.  Burns--Weiss \cite{BurnsWeiss} gives
\[
 \htop(\psi^k)>0,
\]
and hence $\htop(\psi)=\htop(\psi^k)/k>0$.

To pass from the return map to the flow, choose an ergodic
$\psi$-invariant probability measure $\nu$ with $h_\nu(\psi)>0$, which
exists by the variational principle and ergodic decomposition.  The boundary
circle is invariant and the restriction of $\psi$ to it is a circle
diffeomorphism, hence has zero entropy; consequently $\nu$ gives zero mass to
the boundary.  If $\tau$ denotes the positive return time on the interior,
its smooth extension to the closed disk implies
\[
 0<\int\tau\,d\nu<\infty.
\]
The suspension measure determined by $\nu$ is therefore a measure for the
Hamiltonian flow on $\cM\setminus\zeta_e$ and, by Abramov's entropy formula
\cite{Abramov}, has positive entropy.  Therefore
\[
 \htop(\phi_H|_{\cM})>0.
\]
\end{proof}

There is a stronger, more classical alternative.

\begin{proposition}[Transverse version]\label{prop:transverse}
Suppose that, in addition to \eqref{eq:zsign}, the tangent
\[
 \gamma'(u)=\bigl(p_{\rm apo}'(u),z_{\rm apo}'(u)\bigr)
\]
satisfies at every possible zero of $z$
\[
 p_{\rm apo}'(u)\ne0,\qquad z_{\rm apo}'(u)\ne0.
\]
Then the homoclinic intersection is transverse.  In particular a suitable
Poincar\'e map contains a Smale horseshoe.
\end{proposition}

\begin{proof}
At a symmetric homoclinic point write the unstable tangent as $(P',Z')$,
where $P'=p_{\rm apo}'(u)$ and $Z'=z_{\rm apo}'(u)$.  Reflection
\eqref{eq:apoReflection} sends it to the stable tangent $(P',-Z')$.  Their
determinant is
\[
 \det
 \begin{pmatrix}
 P'&P'\\
 Z'&-Z'
 \end{pmatrix}
 =
 -2P'Z'\ne0.
\]
The Smale--Birkhoff homoclinic theorem gives a horseshoe.
\end{proof}

The current numerical search was originally designed for the stronger
targets
\begin{equation}\label{eq:strongtargets}
 P'<-100,\qquad
 Z'>10^5,
\end{equation}
which have large floating-point margins.  For positive entropy alone,
Proposition~\ref{prop:crossing} shows that the $Z'$ estimate may be omitted.

\section{How interval arithmetic is used in the proof}\label{sec:interval}

We now explain the computer-assisted component in more detail.  The purpose
of this section is not to reproduce the implementation line by line, but to
make clear why the numerical statements, once validated, are mathematical
theorems.

\subsection{Intervals and outward rounding}

An interval
\[
 X=[\underline x,\overline x]
\]
represents all real numbers between its endpoints.  Arithmetic operations are
performed with \emph{outward rounding}: if $x\in X$ and $y\in Y$, then the
machine-computed interval $X\circ Y$ is guaranteed to contain every exact
value $x\circ y$ for the corresponding operation $\circ$.

Thus, for example, if a computation gives
\[
 f(X)\subset[0.002,0.014],
\]
then $f(x)>0$ for every $x\in X$.  The statement is not a floating-point
observation; it is a rigorous enclosure statement.

The same idea applies to vectors and matrices.  An interval matrix
$\mathbf A$ contains an entire family of exact matrices.  Inequalities
written using suprema and infima of interval entries therefore hold
simultaneously for every exact derivative matrix represented by
$\mathbf A$.

\subsection{Validated integration of the flow}

The vector field \eqref{eq:HamEq} is explicit and analytic away from
collisions.  We use the CAPD::DynSys library
\cite{KapelaMrozekWilczakZgliczynski}, which provides validated integration
of ordinary differential equations and their variational equations.  For
reproducibility, the accompanying workflow pins the CAPD source to commit
\texttt{731079217a9254ea2948d742df2b170895effe7f}.

For a box $X$ of initial conditions, the validated ODE solver does not follow
a single floating-point trajectory.  Instead it constructs an enclosure
containing
\[
 \phi^t(x)
 \qquad\text{for every }x\in X
\]
throughout each integration step.  In a $C^1$ computation it simultaneously
encloses the solutions of the variational equation
\[
 \dot\Xi(t)=DX_H(\phi^t(x))\Xi(t),
 \qquad \Xi(0)=I.
\]
Consequently it encloses both the flow and its first derivative.

The main practical difficulty is the wrapping effect: after many iterates,
a rectangular enclosure may become much wider than the true image.
CAPD uses doubleton and Hermite--Obreshkov representations to reduce this
effect.  In the present problem the relevant phase-space dimension is only
four, while the Poincar\'e dynamics is two-dimensional, so the required
validation is comparatively small.

\subsection{Validated Poincar\'e maps}

The first-return map is defined by an event:
\[
 \pr=0
\]
with a prescribed crossing direction.  CAPD's interval Poincar\'e map
routines rigorously enclose the first hitting point and the return time.
The variational equation initially gives the derivative of the flow.  The
standard correction for the dependence of the hitting time is then used to
produce an enclosure of the derivative of the Poincar\'e map.

In the implementation we use two event maps:
\begin{itemize}
\item a minus-to-plus crossing of $\pr=0$ for the periapsis map $P$;
\item a plus-to-minus crossing for the periapsis-to-apoapsis map $A$.
\end{itemize}
The sign condition
\[
 \dot{\pr}<-\delta<0
\]
at the validated apoapsis box guarantees that the latter event is regular and
that the hitting map is smooth there.

The Poincar\'e computation takes place on the fixed energy hypersurface.
The lift
\[
 (p,z)\longmapsto(r(p,z),z,0,p)
\]
must therefore also be validated.  Lemma~\ref{lem:energygraph} supplies
existence and uniqueness of $r(p,z)$, while
\eqref{eq:liftderivatives} gives its derivatives.  This is then composed with
the full four-dimensional Poincar\'e derivative to obtain a rigorous
$2\times2$ derivative matrix in $(p,z)$ coordinates.

\subsection{Validating the period-four saddle}

The first scalar validation concerns
\[
 g(p)=\pi_pP(p,0).
\]
The computer proves
\[
 g(8.5768)<0<g(8.5769)
\]
and
\[
 g'(I_*)>0.
\]
The intermediate value theorem gives a unique root $p_*\in I_*$.
A rigorous bisection may then reduce the root interval.

The program evaluates $P^4$ and its derivative on the entire root enclosure.
From the interval trace and determinant it encloses the two eigenvalues and
checks \eqref{eq:multipliertarget}.  This establishes a genuine hyperbolic
period-four orbit; the approximate value $p_*\approx8.5768221605$ is used
only as centering information.

\subsection{Validating the true unstable manifold}

The local coordinate matrix \eqref{eq:B} is also only centering information.
No claim depends on those decimal numbers being exact eigenvectors.

For a candidate block
\[
 N=[-U,U]\times[-aU,aU],
\]
the interval computation encloses $D\widehat F(N)$ and verifies
\eqref{eq:cone1}--\eqref{eq:cone2}.  It then computes interval images of the
whole block and of the two vertical faces to verify
\eqref{eq:trap}--\eqref{eq:exit}.  Proposition~\ref{prop:graph} then converts
these finitely many interval inequalities into a statement about the
\emph{actual} unstable manifold:
\[
 v=h(u),\qquad |h'(u)|\le a.
\]
In particular,
\[
 h(u_j)\in[-a|u_j|,a|u_j|],
 \qquad j=L,R.
\]
These are precisely the two tubes that are propagated in the next step.

\subsection{Validating the homoclinic crossing}

The long map
\[
 G=A\circ P^{24}
\]
is applied to the two endpoint tubes.  If interval arithmetic proves
\[
 \pi_zG(\text{tube at }u_L)<0,
 \qquad
 \pi_zG(\text{tube at }u_R)>0,
\]
then \eqref{eq:zsign} holds for the true unstable graph.

For the crossing condition, one propagates the whole tangent cone
\[
 (1,m),\qquad |m|\le a.
\]
After conversion to physical coordinates and multiplication by the rigorous
derivative of $G$, interval arithmetic encloses the $p$ component of every
possible tangent to the propagated unstable arc.  To control wrapping in the
long map $A\circ P^{24}$, the implementation subdivides
$[u_L,u_R]$ into $32$ subintervals and validates the tangent cone on each one.
It is enough to prove
\[
 P'<-\eta<0.
\]
The present code uses the conservative target $\eta=50$.  This proves
\eqref{eq:pmonotone} without having to know the exact function $h$.

A stronger validation also encloses the $z$ component and proves
\eqref{eq:strongtargets}.  In that case the homoclinic intersection is
transverse.

\subsection{Why this is different from numerical evidence}

The floating-point search serves only to locate promising boxes.  A usual
simulation gives approximate values along one orbit and may miss errors
caused by roundoff, event location, accumulated integration error and
sensitivity to initial data.  The interval proof is instead a finite chain
of inclusions.  Every object used in the mathematical argument---the root,
the multipliers, the image of the block, the endpoint tubes and the tangent
cone---is enclosed with directed rounding.

Therefore a successful run is a proof of a list of explicit inequalities.
The dynamical conclusion follows from the exact propositions of the previous
sections.

\section{The finite certificate}\label{sec:certificate}

We collect the assumptions of Theorem~\ref{thm:certificate} in a form suited
for independent implementation.

\subsection*{(C1) Symmetric hyperbolic period-four point}

The interval computation must prove
\begin{align}
 g(8.5768)&<0<g(8.5769),\label{eq:C1a}\\
 g'([8.5768,8.5769])&>0.\label{eq:C1b}
\end{align}
Let $I_q$ be the resulting root enclosure.  It must also prove that the two
multipliers of $DP^4$ on $I_q\times\{0\}$ satisfy
\begin{equation}\label{eq:C1c}
 \lambda_u\in(5.5,5.8),
 \qquad
 \lambda_s\in(0.17,0.19).
\end{equation}

\subsection*{(C2) Energy lift}

For every local box used in the proof, the lift to $H=-1$ must belong to the
unique branch of Lemma~\ref{lem:energygraph}.  It is sufficient to verify
that the section coordinates remain in
\[
 8.5767\le p\le8.5770,
 \qquad
 |z|\le10^{-4},
\]
and then use \eqref{eq:Hrbound}--\eqref{eq:PhiSides} before applying interval
Newton.  In particular, an implementation that only checks
$0\notin H_r$ after the Newton iteration is not sufficient: it proves
uniqueness on the surviving interval but does not by itself prove existence
of the energy branch for every parameter value in the input box.

\subsection*{(C3) Local hyperbolic block}

For some
\[
 U\in\{4,4.5,5,6,8\}\times10^{-5},
\]
and some aperture $a>0$, with coordinates \eqref{eq:B}, prove
\eqref{eq:cone1}--\eqref{eq:exit}.  Explicitly, a convenient interval form is
\begin{align}
 \inf |J_{11}|-a\sup|J_{12}|&>1,\label{eq:C3a}\\
 \sup|J_{21}|+a\sup|J_{22}|
 &<
 a\bigl(\inf |J_{11}|-a\sup|J_{12}|\bigr),\label{eq:C3b}\\
 \pi_v\widehat F(N)&\Subset(-aU,aU),\label{eq:C3c}\\
 \pi_u\widehat F(\{U\}\times[-aU,aU])&>U,\label{eq:C3d}\\
 \pi_u\widehat F(\{-U\}\times[-aU,aU])&<-U.\label{eq:C3e}
\end{align}

\subsection*{(C4) Homoclinic topological crossing}

For
\[
 u_L=2.68\times10^{-5},
 \qquad
 u_R=2.70\times10^{-5},
\]
propagate the two tubes
\[
 \{u_j\}\times[-a|u_j|,a|u_j|],
 \qquad j=L,R,
\]
under $G=A\circ P^{24}$ and prove
\begin{equation}\label{eq:C4a}
 z_{\rm apo}(u_L)<0<z_{\rm apo}(u_R).
\end{equation}
On a finite subdivision of the complete strip
\[
 [u_L,u_R]\times[-au_R,au_R],
\]
prove, on every substrip, regularity of the apoapsis event,
\begin{equation}\label{eq:C4b}
 \dot{\pr}<-\delta<0,
\end{equation}
and propagation of the tangent cone with
\begin{equation}\label{eq:C4c}
 P'<-\eta<0.
\end{equation}
Any explicit $\delta,\eta>0$ are sufficient.  In the supplied implementation
we take $\delta=0.02$, $\eta=50$, and use $32$ equal subdivisions of the
$u$-interval.

\subsection*{(C4$^+$) Strong transverse certificate}

Instead of only \eqref{eq:C4c}, one may prove the stronger pair
\begin{equation}\label{eq:C4plus}
 P'<-100,\qquad Z'>10^5.
\end{equation}
Then Proposition~\ref{prop:transverse} gives a transverse homoclinic
intersection.

\begin{proof}[Proof of Theorem~\ref{thm:certificate}]
Condition (C1) and Lemma~\ref{lem:period4} give the symmetric period-four
orbit, and \eqref{eq:C1c} makes it hyperbolic.
Condition (C2) makes every local section lift rigorous.
Condition (C3), together with Proposition~\ref{prop:graph}, produces the true
unstable graph and the endpoint tubes containing it.
Condition (C4a) and Proposition~\ref{prop:homoclinic} give a symmetric
homoclinic point.
Conditions (C4b)--(C4c) and Proposition~\ref{prop:crossing} make one such
point a topological crossing.
Theorem~\ref{thm:entropyfromcrossing} then gives positive topological entropy.
If (C4$^+$) holds, Proposition~\ref{prop:transverse} gives a transverse
homoclinic point and hence a Smale horseshoe.
\end{proof}

\section{Exploratory computation and expected margins}\label{sec:exploratory}

For completeness, we record the floating-point values that motivated the
interval boxes.  They are not used as proof.

The symmetric period-four point was found near
\[
 p_*\approx8.5768221605,
\]
with multipliers
\[
 \lambda_u\approx5.6225545,
 \qquad
 \lambda_s\approx0.1778551.
\]
For the center line $v=0$, propagation by $A\circ P^{24}$ gives
approximately
\[
 z_{\rm apo}(2.68\times10^{-5})\approx-1.58\times10^{-2},
\]
and
\[
 z_{\rm apo}(2.70\times10^{-5})\approx1.25\times10^{-2}.
\]
The corresponding finite-difference slopes are approximately
\[
 \frac{dp}{du}\approx-2.08\times10^2,
 \qquad
 \frac{dz}{du}\approx1.42\times10^5,
\]
and the apoapsis event satisfies roughly
\[
 \dot{\pr}\approx-3.10\times10^{-2}.
\]
These numbers explain the conservative target inequalities in
\eqref{eq:multipliertarget}, \eqref{eq:C4b} and \eqref{eq:C4plus}.

The main numerical issue is not proximity to zero, but control of interval
wrapping over the long iterate $P^{24}$.  The supplied certificate therefore
uses $32$ equal subdivisions of $[u_L,u_R]$.  This number is a numerical
choice rather than a mathematical assumption and may be increased if needed.
Since only the sign of $P'$ is required for the Burns--Weiss route, this
subdivision is considerably easier than validating a very large lower bound
for $Z'$.

\section{Discussion}

The argument above isolates a concrete route toward Marco's entropy problem.
Several features seem useful beyond the particular parameter values chosen
here.

First, reversibility is used in an essential way.  It reduces the search for
a two-dimensional stable--unstable intersection to a scalar zero on the
fixed set of an anti-symplectic reverser.  Once an unstable arc is propagated
to the apoapsis section, the equation $z=0$ automatically places the point
on the stable manifold as well.

Second, the local unstable manifold is certified by a cone/block argument
rather than by a high-order polynomial parametrization.  This keeps the
computer-assisted proof short and robust.

Third, for the purpose of positive entropy one does not need to prove
differential transversality.  The sign change of $z$ together with
monotonicity of $p$ forces a topological crossing, and Burns--Weiss applies
even in the presence of higher-order contact.  Differential transversality
is nevertheless desirable because it immediately gives a Smale horseshoe
and persistence under small parameter changes.

Finally, all inequalities in (C1)--(C4) are strict.  If the stronger
transverse certificate is obtained with uniform margins, the same argument
can in principle be repeated with $(\beta,e)$ ranging in a small parameter
box.  This would prove positive entropy on an open set of compact integral
manifolds, providing a stronger partial answer to Marco's question.

\section*{Acknowledgments}

LL is partially supported by the National Natural Science Foundation of
China (Grant number 12401238) and the Natural Science Foundation of Shandong
Province (Grant number ZR2024QA188).  LL thanks the School of Mathematics at
Shandong University for its support.  PS is partially supported by the
National Natural Science Foundation of China (Grant number W2431007).
LL and PS thank the Shenzhen International Center for Mathematics--SUSTech
for its support.

\begin{thebibliography}{99}

\bibitem{Abramov}
L.~M.~Abramov,
\newblock On the entropy of a flow,
\newblock \emph{Dokl. Akad. Nauk SSSR} \textbf{128} (1959), 873--875.

\bibitem{AlbouyCabralSantos}
A.~Albouy, H.~E.~Cabral and A.~A.~Santos,
\newblock Some problems on the classical $n$-body problem,
\newblock \emph{Celestial Mech. Dynam. Astronom.} \textbf{113} (2012),
369--375.
\newblock \href{https://doi.org/10.1007/s10569-012-9431-1}
{doi:10.1007/s10569-012-9431-1}.

\bibitem{BurnsWeiss}
K.~Burns and H.~Weiss,
\newblock A geometric criterion for positive topological entropy,
\newblock \emph{Comm. Math. Phys.} \textbf{172} (1995), 95--118.
\newblock \href{https://doi.org/10.1007/BF02104512}
{doi:10.1007/BF02104512}.

\bibitem{HuLiuOuSalomaoYu}
X.~Hu, L.~Liu, Y.~Ou, P.~A.~S.~Salom\~ao and G.~Yu,
\newblock A symplectic dynamics approach to the spatial isosceles three-body
problem,
\newblock \emph{J. Eur. Math. Soc.} (2025), published online first.
\newblock \href{https://doi.org/10.4171/JEMS/1678}{doi:10.4171/JEMS/1678}.

\bibitem{ECHconstraints}
X.~Hu, L.~Liu, Y.~Ou, Z.~Qiao and P.~A.~S.~Salom\~ao,
\newblock ECH constraints and twist dynamics in the spatial isosceles
three-body problem,
\newblock preprint, arXiv:2602.24025.

\bibitem{KapelaMrozekWilczakZgliczynski}
T.~Kapela, M.~Mrozek, D.~Wilczak and P.~Zgliczy\'nski,
\newblock CAPD::DynSys: a flexible C++ toolbox for rigorous numerical
analysis of dynamical systems,
\newblock \emph{Commun. Nonlinear Sci. Numer. Simul.} \textbf{101} (2021),
105578.
\newblock \href{https://doi.org/10.1016/j.cnsns.2020.105578}
{doi:10.1016/j.cnsns.2020.105578}.

\bibitem{Moeckel1984}
R.~Moeckel,
\newblock Heteroclinic phenomena in the isosceles three-body problem,
\newblock \emph{SIAM J. Math. Anal.} \textbf{15} (1984), 857--876.
\newblock \href{https://doi.org/10.1137/0515065}
{doi:10.1137/0515065}.

\end{thebibliography}

\end{document}
